% ===========================================================================
\chapter{Turn by turn}\label{ch:directions}

A route is a list of nodes: from the beach car park to the lighthouse, a
couple of dozen of them. Nobody wants to hear about nodes. They want to
hear ``Head west on Coast Road. After 3.8 km, turn left onto Lighthouse
Lane.'' This chapter turns routes into instructions, and then into the
first complete program of the project.

\section{Where to say something}

Most nodes of a route need no instruction at all: the traveller just keeps
going along the same road. Something needs saying only where the road
changes. So we walk along the route, edge by edge, and cut it into
\emph{stretches}, each along a single road. Every stretch begins with an
instruction:

\begin{itemize}[nosep]
  \item the first stretch with which way to set off: ``Head north on
    Harbour Street'';
  \item every later stretch with the turn onto its road: ``Turn left onto
    Mill Lane'';
\end{itemize}
and a last instruction says that the traveller has arrived.

How sharp is a turn? At the node where the road changes, compare the
bearing of the segment arriving there with the bearing of the segment
leaving it; \code{turn_angle} from Chapter~\ref{ch:geometry} gives the
difference, positive to the right. Then a few thresholds put it into words:

\begin{center}
\begin{tikzpicture}[scale=1.0]
  \draw[color=cm, thick, -{Stealth}] (0,-1.6) -- (0,0);
  \foreach \a/\lab in {0/Continue, 40/Bear right, -40/Bear left, 95/Turn right,
      -95/Turn left, 160/Turn sharp right, -160/Turn sharp left} {
    \draw[accent, -{Stealth}] (0,0) -- ({1.6*sin(\a)},{1.6*cos(\a)});
    \node[font=\scriptsize, text=accent] at ({2.35*sin(\a)},{2.1*cos(\a)}) {\lab};
  }
  \fill (0,0) circle (0.05);
  \node[font=\scriptsize, text=cm] at (0.6,-1.2) {arriving};
\end{tikzpicture}
\end{center}

Less than $25^\circ$ either way is carrying straight on; up to $60^\circ$
is bearing left or right; up to $140^\circ$ is turning; more than that is
turning sharply, nearly back the way you came.

\begin{newfile}{directions.nex}
-- directions.nex: turning a route into instructions a person can follow.

intern Search
intern Helpers

-- One instruction: what to do, where, and how far into the route.
class Instruction
create
  make(words: String, at: Integer, metres, degrees: Real) do
    text := words
    node := at
    distance := metres
    turn := degrees
  end
feature
  -- What to do: "Turn left onto Mill Lane".
  once text: String
  -- The node where it happens.
  once node: Integer
  -- How far it is from the start of the route, in metres.
  once distance: Real
  -- How sharply the traveller turns there, in degrees: positive to the
  -- right, negative to the left, 0 for straight on.
  once turn: Real
end

-- The words for a turn of `angle` degrees (positive to the right).
function manoeuvre(angle: Real): String do
  let side: String := when angle > 0.0 then "right" else "left" end
  if angle.abs < 25.0 then
    result := "Continue"
  elseif angle.abs < 60.0 then
    result := "Bear " + side
  elseif angle.abs < 140.0 then
    result := "Turn " + side
  else
    result := "Turn sharp " + side
  end
end
\end{newfile}

Each instruction remembers how far into the route it is. The navigator
will need that, to say ``in 300 m''. It also keeps the turn's angle, so
that the live map can draw an arrow bending the right way.

\section{The instructions}

\begin{append}{directions.nex}
-- Instructions for following `path` to `destination`.
--
-- A route can have dozens of nodes, but a person needs to be told
-- something only where the road changes. So the route is cut into
-- stretches along one road each. The first instruction says which way to
-- set off; each later stretch starts with the turn onto its road; and the
-- last instruction says that the traveller has arrived.
function directions(g: Road_Graph, path: Array[Integer], destination: String): Array[Instruction]
require
  is_a_route: is_route(g, path)
do
  let pts: Array[Point] := g.island.points_of(path)
  if path.length = 1 then
    result.add(create Instruction.make("You are at " + destination, path.get(0), 0.0, 0.0))
  else
    let travelled: Real := 0.0
    let previous: ?Edge := nil
    from
      let i: Integer := 0
    until
      i = path.length - 1
    do
      if ?g.edge_between(path.get(i), path.get(i + 1)) as e then
        if ?previous as before then
          if e.road.name /= before.road.name then
            let angle: Real := turn_angle(bearing(pts.get(i - 1), pts.get(i)),
              bearing(pts.get(i), pts.get(i + 1)))
            result.add(create Instruction.make(manoeuvre(angle) + " onto " + e.road.name,
              path.get(i), travelled, angle))
          end
        else
          result.add(create Instruction.make("Head "
            + compass_name(bearing(pts.get(0), pts.get(1))) + " on " + e.road.name,
            path.get(0), 0.0, 0.0))
        end
        travelled := travelled + e.length
        previous := e
      end
      i := i + 1
    end
    result.add(create Instruction.make("Arrive at " + destination,
      path.get(path.length - 1), travelled, 0.0))
  end
ensure
  starts_at_start: result.get(0).distance = 0.0
  ends_at_end: (result.get(result.length - 1).distance
    - route_cost(g, path, by_distance())).abs < 0.001
end
\end{append}

\code{previous} is the edge before the current one, or \code{nil} for the
very first edge. \code{?previous as before} is Nex's \emph{object test}:
it is true when \code{previous} is not \code{nil}, and inside the
\code{if} it names the edge \code{before}, with the type \code{Edge}
instead of \code{?Edge}, so that we can use its road. On the first edge
there is no edge before, so that is where we say which way to set off.

The postcondition checks that the distances add up: the last instruction
is exactly as far into the route as the route is long.

Real navigators say a good deal more than ours: ``keep left'' where a road
forks, ``take the second exit'' at a roundabout, ``turn right after the
church'' when there is a landmark, and they warn about turns where the road
keeps its name but the route leaves it. They all start from the same idea:
find the places where something changes, and describe the change.

\section{A route planner}

Now everything comes together in a program. It takes two place names, and
a few options, from the command line, plans a route, and prints its
instructions, with the distance of each stretch.

\begin{newfile}{route.nex}
-- route.nex: directions from one place to another, printed.
--
--   nex route.nex                              -- list the places
--   nex route.nex "The Harbour" Lighthouse     -- the quickest way by car
--   nex route.nex harbour lighthouse --walk    -- on foot
--   nex route.nex harbour lighthouse --shortest
--   nex route.nex harbour lighthouse --algorithm dijkstra   (or bfs, astar)
--
-- A place can be given by part of its name: "mill" is "Old Mill".

intern Directions
intern Helpers

function main() do
  let args: Array[String] := program.arguments
  let names: Array[String] := []
  let mode: Mode := Mode.Driving
  let shortest: Boolean := false
  let algorithm: String := "astar"
  let map_file: String := "kestrel_island.map"
  from
    let i: Integer := 0
  until
    i >= args.length
  do
    let arg: String := args.get(i)
    if arg = "--walk" then
      mode := Mode.Walking
    elseif arg = "--shortest" then
      shortest := true
    elseif arg = "--algorithm" and i + 1 < args.length then
      algorithm := args.get(i + 1)
      i := i + 1
    elseif arg = "--map" and i + 1 < args.length then
      map_file := args.get(i + 1)
      i := i + 1
    else
      names.add(arg)
    end
    i := i + 1
  end
\end{newfile}

\code{program.arguments} gives the words typed after the
program's name. Options start with \code{--}; anything else is a place
name. Next, with the map loaded, either explain how to use the program,
or plan the route:

\begin{continued}{route.nex}
  let island: Island_Map := create Island_Map.load(user_path(map_file))
  if names.length /= 2 then
    console.print_line("usage: nex route.nex FROM TO [--walk] [--shortest] [--algorithm bfs|dijkstra|astar]")
    console.print_line("")
    console.print_line("Places on " + island.name + ":")
    across island.places as p do
      console.print_line("  " + p.name)
    end
  elseif ?island.find_place(names.get(0)) as start
      and ?island.find_place(names.get(1)) as finish then
    let g: Road_Graph := create Road_Graph.build(island, mode)
    let source: Integer := g.nearest_node(island.point(start.node_id))
    let target: Integer := g.nearest_node(island.point(finish.node_id))
    let found: Search_Result := find_route(g, source, target, algorithm, shortest)
    if found.found() then
      print_route(g, found, start.name, finish.name)
    else
      console.print_line("There is no way from " + start.name + " to " + finish.name + ".")
    end
  else
    console.print_line("I don't know one of those places. Run nex route.nex to list them.")
  end
end
\end{continued}

The instructions are printed numbered, each with the length of its stretch:
the distance from it to the next instruction.

\begin{continued}{route.nex}
function print_route(g: Road_Graph, found: Search_Result, from_name, to_name: String) do
  let how: String := when g.mode = Mode.Walking then "Walking" else "Driving" end
  console.print_line(how + " from " + from_name + " to " + to_name + ": "
    + distance_text(route_cost(g, found.path, by_distance())) + ", about "
    + duration_text(route_cost(g, found.path, by_time())) + ".")
  console.print_line(found.algorithm + " explored " + found.explored.length.to_string
    + " of the " + g.nodes.length.to_string + " nodes.")
  console.print_line("")
  let steps: Array[Instruction] := directions(g, found.path, to_name)
  from
    let i: Integer := 0
  until
    i = steps.length
  do
    let step: Instruction := steps.get(i)
    let line: String := (i + 1).to_string.pad_start(" ", 3) + ". " + step.text
    if i + 1 < steps.length then
      line := line.pad_end(" ", 42)
        + distance_text(steps.get(i + 1).distance - step.distance).pad_start(" ", 8)
    end
    console.print_line(line)
    i := i + 1
  end
end

main()
\end{continued}

\section{Planning some trips}

With no places, the program lists them:

\runoutput{route.nex}{route_places}

From the beach to the lighthouse, by car, the quickest way and the
shortest:

\runoutput{route.nex "beach car park" lighthouse}{route_beach}
\runoutput{route.nex beach lighthouse --shortest}{route_beach_shortest}

From Market Square to the harbour, a few hundred metres, by car and on
foot. The one-way streets send the car east, then south down Chapel
Street, and back west along the quay; the walker goes straight down
Harbour Street:

\runoutput{route.nex market harbour}{route_market_drive}
\runoutput{route.nex market harbour --walk}{route_market_walk}

And from Kestrel Peak, in the middle of Hollow Wood, to the lighthouse. A
car must take the track out of the wood and go round by the Coast Road;
on foot, the paths over the hill and along the cliffs are much shorter:

\runoutput{route.nex peak lighthouse}{route_peak_drive}
\runoutput{route.nex peak lighthouse --walk}{route_peak_walk}

\section{Checking the directions}

Two more checks: that the directions start and end the way they should, and
that the one-way streets really do send drivers round by Chapel Street
while walkers go the short way.

\begin{addbefore}{checks.nex}{before main}{function main() do}
function check_directions(c: Checker, m: Island_Map, driving, walking: Road_Graph) do
  if ?m.find_place("market") as a and ?m.find_place("harbour") as b then
    let drive: Search_Result := dijkstra(driving, node_of(driving, a), node_of(driving, b), by_time())
    let steps: Array[Instruction] := directions(driving, drive.path, b.name)
    c.check("directions start by setting off", steps.get(0).text.starts_with("Head "))
    c.check("directions end by arriving", steps.get(steps.length - 1).text = "Arrive at The Harbour")
    let detour: Boolean := false
    across steps as step do
      if step.text = "Turn right onto Chapel Street" then
        detour := true
      end
    end
    c.check("one-way streets send drivers round by Chapel Street", detour)
    let walk: Search_Result := dijkstra(walking, node_of(walking, a), node_of(walking, b), by_time())
    c.check("walkers go the short way", route_cost(walking, walk.path, by_distance())
      < route_cost(driving, drive.path, by_distance()))
  end
  c.check("manoeuvres", manoeuvre(5.0) = "Continue" and manoeuvre(-40.0) = "Bear left"
    and manoeuvre(90.0) = "Turn right" and manoeuvre(-170.0) = "Turn sharp left")
end
\end{addbefore}

\begin{tightaddbefore}{checks.nex}{to main, before c.summary()}{  c.summary()}
  check_directions(c, m, driving, walking)
\end{tightaddbefore}

The checks now use \code{directions.nex}, so we intern it at the top of the
file. (\code{directions.nex} interns \code{search.nex} itself, so the
second \code{intern} is no longer needed, but interning a file twice is
harmless.)

\begin{tightaddbefore}{checks.nex}{at the top}{intern Search}
intern Directions
\end{tightaddbefore}
